\begin{abstract}
\noindent I estimate a drift-diffusion model of snack choices in Experiment 1,
using differences in bids to measure relative value. My 435 choices yield a
choice-sensitivity estimate of 0.408, compared with 0.308 for the 30,450 choices
pooled across the class. The fitted probabilities rise with the bid difference
in both samples. Pooled choices change more sharply near equal bids and remain
less extreme in the tails than the fitted curve predicts. Response times
provide a second check because their shape depends on the same choice
parameter. Holding that estimate fixed, I fit the time scale and nondecision
time to 15 bin averages. The model explains 55.4 percent of the variation
across my trimmed response-time means and 96.1 percent across the class means.
The class averages follow the predicted inverted-U closely. My data show a
less regular pattern, including faster choices near equal bids than in some
neighboring bins.
\end{abstract}

\vspace{1em}
\noindent \textbf{JEL Codes:} D81, D87

\vspace{0.25em}
\noindent \textbf{Keywords:} drift-diffusion model, stochastic choice, response time
